\begin{abstract}
Developers choose open-source licenses, or inherit them through dependencies, often without understanding the obligations involved. Large language models (LLMs) explain legal text fluently but also state unsupported facts, which is costly in licensing. \sys{} is a two-stage LLM\,+\,NLI system: an open-weight LLM recommends three licenses from a catalogue of 34 OSI-approved software licenses and explains each with plain-language claims that quote a specific clause, and an independent natural language inference (NLI) model marks every claim \emph{verified}, \emph{unverified}, or \emph{contradicted}. All ground truth is computed from public, curated license data, so no legal expert is needed. Phase~1 showed two weaknesses: the recommender broke stated requirements (top-pick accuracy 0.56), and the off-the-shelf verifier confirmed only 69\% of true claims. In Phase~2 we (1) grounded the recommender in the catalogue's license labels with a check-then-rank prompt, raising top-pick accuracy from 0.558 to 0.885 and top-3 accuracy from 0.740 to 0.942 on the same 104 descriptions, while a smaller model with the old prompt did slightly worse; (2) fine-tuned the top four layers of DeBERTa-v3-large on 729 clause--claim pairs, raising verifier accuracy from 0.815 to 0.925 and \emph{verified} precision from 0.944 to 0.978 on 146 human-labelled pairs (0.911 and 0.933 when the claim wording is unseen); and (3) turned the prototype into a connected web application with a FastAPI backend, a SQLite job store, a 27-test suite, and GitHub Actions CI on a two-branch workflow. Requests for weak copyleft remain the main failure case.
\end{abstract}

\begin{IEEEkeywords}
open-source licensing, large language models, natural language inference, explainability, fine-tuning, software engineering for ML
\end{IEEEkeywords}
